\section{Conclusion}
Healthcare agents should be evaluated according to what they observe, propose, commit, and verify, together with the authority that permits those effects. The selected evidence spans informational, interactive, artifact, and sandbox capabilities; it does not justify inferring live clinical benefit from aggregate benchmark success. Clinical Action Assurance provides a proposed structure for patient binding, evidence provenance, approval-bound contracts, current-state admission, bounded execution, and monitored effect verification. Its utility and clinical tradeoffs require further empirical assessment. Progress toward reliable clinical agency depends on connecting execution-level evidence with human workload, omissions, equitable access, and patient outcomes.
